\section{Further research questions}
\label{sec:further-research}

The proof architecture suggests several concrete follow-up projects.

\subsection{Finish the object-theory bridge}

The highest-priority task is to implement the module tree of
Section~\ref{sec:lean-plan} and produce the first end-to-end theorem
\[
 []\vdash_{\ipra}
 \forall c\;\neg\operatorname{UniversalRadicalFormula}(5,c).
\]
A useful engineering question is whether the ordinary \code{PRAAx} proof or a
direct proof in one of ConciseIPRA's smaller clause calculi yields the smaller
checked certificate.  The bi-interpretation guarantees equivalence, but proof
sizes may differ substantially.

\subsection{Replace automorphism tables by permutation presentations}

The radical-normal-closure certificate proposed here can list finite
multiplication tables.  Those tables may be enormous.  A more economical
certificate would present each group as a permutation group generated by a
short list and prove solvability through a subnormal series with abelian factor
certificates.  All membership tests remain bounded and primitive recursive.
The research problem is to find the smallest checker and certificate format
whose soundness still has a transparent iPRA proof.

\subsection{A branch-free constructive quartic object}

The report uses the decidable coefficient guard $q=0$ versus $q\ne0$.
One can instead build a universal finite splitting algebra in which the
Ferrari parameter equations are imposed simultaneously, then prove a product
identity without selecting a branch.  When is that algebra guaranteed to be a
nonzero finite product of fields constructively?  A satisfactory answer could
yield a formula theorem over fields without decidable equality while retaining
an explicit computational interpretation.

\subsection{A rational counterexample in iPRA}

Section~\ref{sec:explicit-quintic} reduces the rational example
$X^5-X-1$ to finite modular arithmetic plus the Dedekind--Frobenius cycle
lemma.  Formalizing that lemma in iPRA would give a closed integer-coefficient
counterexample, avoiding rational-function fields in the final sentence.
ProveIt's existing Selmer development supplies a strong semantic target.  A
particularly attractive endpoint is a numeral-only certificate proving that
no root of the displayed quintic has a radical-expression code.

\subsection{Proof-size lower and upper bounds}

The mathematical proof is elementary, but its naively expanded arithmetic
certificate may be very large.  Natural quantitative questions are:

\begin{itemize}
\item What is the smallest sparse-ring trace for the Cardano and Ferrari
  identities?
\item Is the explicit perfectness proof of $A_5$ smaller than evaluation of a
  generated subgroup table after compression?
\item How much does a subnormal-series certificate reduce the radical-tower
  proof compared with a derived-series certificate?
\item What is the compressed size of the final proof under the current
  ConciseIPRA checker metric?
\end{itemize}

These are mathematical proof-compression questions, not merely implementation
tuning: a better certificate format can expose a simpler invariant and lead to
a genuinely shorter iPRA derivation.

\subsection{Solvability decisions beyond degree five}

ProveIt already contains primitive-recursive or recursive decision procedures
for individual quintics and sextics.  The present generic obstruction and
those decision procedures could be connected at the code level: the same
\code{RadicalExpression} semantics should be reflected both by a positive
solver certificate and by a negative Galois-group certificate.  A unified
interface would make it possible to return one of two proof objects for an
input polynomial: a complete radical factorization or a checked nonsolvable
Galois-group witness.
